\documentclass[a4paper,12pt]{article}

% -------------------------------------------------
% Кодировка и язык
% -------------------------------------------------
\usepackage[T2A]{fontenc}
\usepackage[utf8]{inputenc}
\usepackage[russian]{babel}

% -------------------------------------------------
% Математика
% -------------------------------------------------
\usepackage{amsmath,amssymb,amsthm,mathtools}
\usepackage{icomma}

% -------------------------------------------------
% Шрифт и типографика
% -------------------------------------------------
\usepackage{helvet}
\renewcommand{\familydefault}{\sfdefault}
\usepackage{microtype}
\usepackage{setspace}
\onehalfspacing
\usepackage{parskip}

% -------------------------------------------------
% Страница
% -------------------------------------------------
\usepackage[
    left=20mm,
    right=20mm,
    top=20mm,
    bottom=22mm,
    headheight=15pt
]{geometry}

% -------------------------------------------------
% Графика
% -------------------------------------------------
\usepackage{tikz}
\usetikzlibrary{
    calc,
    arrows.meta,
    positioning,
    shapes.geometric
}
\usepackage{tkz-euclide}
\usepackage{pgfplots}
\pgfplotsset{compat=1.18}

% -------------------------------------------------
% Таблицы и списки
% -------------------------------------------------
\usepackage{tabularx,booktabs,longtable}
\usepackage{enumitem}

% -------------------------------------------------
% Служебные пакеты
% -------------------------------------------------
\usepackage{xcolor}
\usepackage{hyperref}
\usepackage{fancyhdr}

% -------------------------------------------------
% Цветовая палитра
% -------------------------------------------------
\definecolor{accent}{HTML}{008080}
\definecolor{darkaccent}{HTML}{004D4D}
\definecolor{textgray}{HTML}{333333}
\definecolor{softgray}{HTML}{6F7777}
\definecolor{lightaccent}{HTML}{DCEEEE}

\color{textgray}

% -------------------------------------------------
% Гиперссылки
% -------------------------------------------------
\hypersetup{
    colorlinks=true,
    linkcolor=darkaccent,
    urlcolor=darkaccent
}

% -------------------------------------------------
% Колонтитулы
% -------------------------------------------------
\pagestyle{fancy}
\fancyhf{}
\fancyhead[L]{\small\color{softgray}Степени и показательная функция}
\fancyhead[R]{\small\color{softgray}ЕГЭ}
\fancyfoot[C]{\small\color{softgray}\thepage}
\renewcommand{\headrulewidth}{0pt}

% -------------------------------------------------
% Заголовки
% -------------------------------------------------
\usepackage{titlesec}

\titleformat{\section}
  {\Large\bfseries\color{darkaccent}}
  {}
  {0pt}
  {}

\titleformat{\subsection}
  {\large\bfseries\color{accent}}
  {}
  {0pt}
  {}

\titlespacing*{\section}{0pt}{1.4em}{0.65em}
\titlespacing*{\subsection}{0pt}{1em}{0.4em}

% -------------------------------------------------
% Настройка списков
% -------------------------------------------------
\setlist[itemize]{
    leftmargin=1.5em,
    itemsep=0.25em,
    topsep=0.35em
}

\setlist[enumerate]{
    leftmargin=1.8em,
    itemsep=0.55em,
    topsep=0.4em
}

% -------------------------------------------------
% TikZ: блок-схемы
% -------------------------------------------------
\tikzset{
    flowstart/.style={
        ellipse,
        draw=accent,
        line width=0.9pt,
        minimum width=39mm,
        minimum height=11mm,
        align=center,
        font=\small
    },
    flowaction/.style={
        rounded corners=3pt,
        rectangle,
        draw=accent,
        line width=0.9pt,
        text width=62mm,
        minimum height=14mm,
        align=center,
        inner sep=4pt,
        font=\small
    },
    flowdecision/.style={
        diamond,
        aspect=2.2,
        draw=accent,
        line width=0.9pt,
        text width=38mm,
        align=center,
        inner sep=1.5pt,
        font=\small
    },
    flowarrow/.style={
        -{Latex[length=3mm]},
        line width=0.9pt,
        draw=darkaccent
    }
}

% -------------------------------------------------
% Вспомогательные команды
% -------------------------------------------------
\newcommand{\important}[1]{\textbf{\color{darkaccent}#1}}
\newcommand{\keyidea}[1]{%
    \par\medskip
    {\color{accent}\rule{\linewidth}{0.6pt}}\par
    \smallskip
    #1
    \smallskip
    {\color{accent}\rule{\linewidth}{0.6pt}}\par
    \medskip
}

\begin{document}

% =================================================
% ТИТУЛЬНЫЙ ЛИСТ
% =================================================
\begin{titlepage}
\thispagestyle{empty}

\begin{center}

\vspace*{24mm}

{\Large\bfseries\color{darkaccent}
ЧЕК-ЛИСТ
\par}

\vspace{10mm}

{\huge\bfseries
Степени и показательная функция
\par}

\vspace{5mm}

{\Large
вычисления, прикладные модели и графики
\par}

\vspace{22mm}

{\large
ЕГЭ по профильной математике
\par}

\vspace{7mm}

{\large
\today
\par}

\vfill

{\large
Лёвин Артём Александрович
\par}

\vspace{3mm}

{\large\bfseries
эксклюзивно для Екатерины Гнедковой
\par}

\vspace{25mm}

\begin{tikzpicture}[remember picture,overlay]
    \draw[
        accent,
        line width=1.3pt
    ]
    ([xshift=20mm,yshift=22mm]current page.south west)
        .. controls
        ([xshift=70mm,yshift=34mm]current page.south west)
        and
        ([xshift=-70mm,yshift=14mm]current page.south east)
        ..
    ([xshift=-20mm,yshift=25mm]current page.south east);
\end{tikzpicture}

\end{center}
\end{titlepage}

% =================================================
% ОСНОВНОЙ ЧЕК-ЛИСТ
% =================================================

\section{1. Степени: базовые преобразования}

\subsection{Главные свойства}

Для допустимых значений переменных:

\[
a^m\cdot a^n=a^{m+n},
\]

\[
\frac{a^m}{a^n}=a^{m-n},
\qquad a\ne0,
\]

\[
(a^m)^n=a^{mn},
\]

\[
(ab)^n=a^n b^n,
\]

\[
\left(\frac{a}{b}\right)^n
=
\frac{a^n}{b^n},
\qquad b\ne0,
\]

\[
a^{-n}=\frac{1}{a^n},
\qquad a\ne0,
\]

\[
a^0=1,
\qquad a\ne0.
\]

\keyidea{
\important{Степени показателей складываются или вычитаются только при умножении или делении степеней с одинаковым основанием.}

Например,
\[
\frac{3^7}{3^2}=3^5.
\]

Из выражения
\[
3^7-3^2
\]
показатели вычитать нельзя.
}

\subsection{Степень суммы и разности}

Следует особенно внимательно различать:

\[
a^{m-n}
=
\frac{a^m}{a^n},
\]

но

\[
a^{m-n}\ne a^m-a^n.
\]

Аналогично:

\[
a^{m+n}=a^m a^n,
\]

но

\[
a^{m+n}\ne a^m+a^n.
\]

\paragraph{Пример.}
\[
2^{2x-2}
=
\frac{2^{2x}}{2^2}.
\]

Запись
\[
2^{2x-2}=2^{2x}-2^2
\]
ошибочна.

\subsection{Число \(1\) как степень}

При \(a>0\)

\[
1=a^0.
\]

Этот приём особенно полезен при решении показательных уравнений.

Например,

\[
1=5^{x-3}
\]

можно записать как

\[
5^0=5^{x-3},
\]

откуда

\[
x-3=0,
\qquad
x=3.
\]

\subsection{Десятичные дроби полезно переводить в степени}

Например,

\[
0{,}5=\frac12=2^{-1},
\]

\[
0{,}25=\frac14=2^{-2},
\]

\[
0{,}125=\frac18=2^{-3}.
\]

Такой переход часто позволяет сразу привести обе части показательного уравнения к одному основанию.

\paragraph{Пример.}
\[
2^{x-4}=0{,}5.
\]

Так как

\[
0{,}5=2^{-1},
\]

получаем

\[
2^{x-4}=2^{-1},
\]

следовательно,

\[
x-4=-1,
\qquad
x=3.
\]

% =================================================
\section{2. Стандартный вид числа}

Число в стандартном виде записывается как

\[
a\cdot10^n,
\qquad
1\le a<10.
\]

\subsection{Сдвиг десятичной запятой}

Если запятую в коэффициенте сдвигаем \important{вправо на один разряд}, показатель степени \(10\) уменьшается на \(1\):

\[
3{,}2\cdot10^{-8}
=
32\cdot10^{-9}.
\]

Если запятую сдвигаем \important{влево на один разряд}, показатель степени увеличивается на \(1\):

\[
32\cdot10^7
=
3{,}2\cdot10^8.
\]

\keyidea{
При изменении положения запятой коэффициент и степень десяти должны компенсировать друг друга, чтобы значение числа сохранилось.
}

\paragraph{Самопроверка.}
Числа

\[
3{,}2\cdot10^{-8}
\quad\text{и}\quad
32\cdot10^{-9}
\]

равны, поскольку увеличение коэффициента в \(10\) раз компенсируется уменьшением степени десятки на \(1\).

% =================================================
\section{3. Прикладные задачи со степенями}

В задачах ЕГЭ формула часто уже дана. Основная работа состоит в том, чтобы:

\begin{enumerate}
    \item правильно перевести условие на язык равенств;
    \item аккуратно подставить данные;
    \item применить свойства степеней;
    \item проверить арифметику и ограничения.
\end{enumerate}

\subsection{«Уменьшился в \(k\) раз»}

Если величина \(V\) уменьшилась в \(3\) раза, то

\[
V_2=\frac{V_1}{3}.
\]

Эквивалентная запись:

\[
V_1=3V_2.
\]

Вторая форма иногда удобнее, поскольку позволяет избежать дробей.

\subsection{«Увеличился в \(k\) раз»}

Если давление увеличилось в \(9\) раз:

\[
p_2=9p_1.
\]

\subsection{Пример с параметром степени}

Пусть для двух состояний выполняется закон

\[
p_1V_1^\alpha=p_2V_2^\alpha.
\]

Известно:

\[
V_1=3V_2,
\qquad
p_2=9p_1.
\]

Подставим:

\[
p_1(3V_2)^\alpha
=
9p_1V_2^\alpha.
\]

Раскрываем степень произведения:

\[
p_1\cdot3^\alpha V_2^\alpha
=
9p_1V_2^\alpha.
\]

При \(p_1\ne0\) и \(V_2\ne0\) сокращаем одинаковые множители:

\[
3^\alpha=9.
\]

Представляем \(9\) как степень тройки:

\[
3^\alpha=3^2.
\]

Следовательно,

\[
\alpha=2.
\]

\keyidea{
Сокращать можно только \important{общие множители}.

Например, в дроби
\[
\frac{9}{3^\alpha}
\]
нельзя ``сократить \(9\) и \(3\)'' как обычные коэффициенты. Полезно сначала записать
\[
9=3^2,
\]
после чего
\[
\frac{3^2}{3^\alpha}=3^{2-\alpha}.
\]
}

% =================================================
\section{4. Показательная функция}

Основная модель занятия:

\[
f(x)=a^x+b,
\]

где

\[
a>0,
\qquad
a\ne1.
\]

Число \(a\) задаёт основание показательной функции, \(b\) отвечает за вертикальный сдвиг графика.

\subsection{Что означает \(f(c)\)}

Запись

\[
f(6)
\]

означает: подставить

\[
x=6
\]

в формулу функции и вычислить результат.

Например, если

\[
f(x)=2^x-3,
\]

то

\[
f(6)=2^6-3=64-3=61.
\]

\subsection{Что означает \(f(x)=c\)}

Если требуется найти \(x\), при котором

\[
f(x)=3,
\]

необходимо подставить \(3\) вместо значения функции и решить полученное уравнение.

Например, для

\[
f(x)=2^x-5
\]

условие

\[
f(x)=3
\]

даёт

\[
2^x-5=3,
\]

\[
2^x=8,
\]

\[
x=3.
\]

% =================================================
\section{5. Метод узловых точек}

\subsection{Что такое узловая точка}

Узловой будем называть точку графика, у которой обе координаты являются удобными точными числами, прежде всего целыми.

Если точка

\[
A(x_0;y_0)
\]

лежит на графике

\[
y=f(x),
\]

то обязательно выполняется

\[
y_0=f(x_0).
\]

Именно это позволяет получить уравнение для неизвестных коэффициентов.

\subsection{Сколько точек требуется}

Количество независимых точек должно быть достаточным для нахождения неизвестных коэффициентов.

Например, в функции

\[
f(x)=a^x+b
\]

имеются два неизвестных параметра:

\[
a,\qquad b.
\]

Поэтому обычно нужны \important{две информативные точки}.

\subsection{Какие точки выбирать}

В первую очередь удобно рассматривать точки с небольшими координатами:

\[
x=0,\quad x=1,\quad x=-1.
\]

Причины:

\[
a^0=1,
\qquad
a^1=a,
\qquad
a^{-1}=\frac1a.
\]

Однако удобная точка должна давать достаточно информации.

% =================================================
\subsection{Пример: \(f(x)=a^x+b\)}

Пусть график проходит через точки

\[
D(0;-2),
\qquad
C(1;-1).
\]

Подставляем координаты точки \(D\):

\[
-2=a^0+b,
\]

\[
-2=1+b,
\]

откуда

\[
b=-3.
\]

Используем точку \(C\):

\[
-1=a^1+b.
\]

Подставляем \(b=-3\):

\[
-1=a-3,
\]

следовательно,

\[
a=2.
\]

Итак,

\[
\boxed{f(x)=2^x-3}.
\]

Например,

\[
f(6)=2^6-3=61.
\]

\begin{center}
\begin{tikzpicture}
\begin{axis}[
    width=0.76\linewidth,
    height=7.2cm,
    axis lines=middle,
    xmin=-2.2,
    xmax=4.2,
    ymin=-4.2,
    ymax=6.3,
    xtick={-2,-1,0,1,2,3,4},
    ytick={-4,-3,-2,-1,0,1,2,3,4,5,6},
    xlabel={$x$},
    ylabel={$y$},
    samples=250,
    domain=-2:3.2,
    clip=false,
    unit vector ratio=1 1,
    grid=both,
    major grid style={line width=.25pt,draw=gray!28},
    minor tick num=0
]
\addplot[
    accent,
    line width=1.3pt
] {2^x-3};

\addplot[
    only marks,
    mark=*,
    mark size=2.4pt,
    darkaccent
] coordinates {(0,-2) (1,-1)};

\node[anchor=south east] at (axis cs:0,-2) {$D(0;-2)$};
\node[anchor=south west] at (axis cs:1,-1) {$C(1;-1)$};

\end{axis}
\end{tikzpicture}
\end{center}

% =================================================
\section{6. Когда удобная точка оказывается бесполезной}

Рассмотрим функцию

\[
f(x)=a^x.
\]

Каждый её график проходит через точку

\[
(0;1),
\]

поскольку

\[
a^0=1.
\]

Если подставить эту точку:

\[
1=a^0,
\]

получим

\[
1=1.
\]

Это верное равенство, но оно не позволяет найти \(a\).

Следовательно, требуется другая точка.

Например, если график проходит через

\[
(-1;2),
\]

то

\[
2=a^{-1}.
\]

То есть

\[
2=\frac1a,
\]

откуда

\[
a=\frac12.
\]

Формула функции:

\[
f(x)=\left(\frac12\right)^x
=
2^{-x}.
\]

\keyidea{
Точка с нулевой координатой часто удобна, однако после подстановки необходимо проверить, сохранились ли в уравнении неизвестные параметры.
}

% =================================================
\clearpage
\thispagestyle{plain}

\begin{center}
{\Large\bfseries\color{darkaccent}
Алгоритм: восстановление формулы по графику
}
\end{center}

\vspace{7mm}

\begin{center}
\begin{tikzpicture}[
    node distance=10mm and 18mm
]

\node[flowstart] (start)
{Дан график функции\\с неизвестными коэффициентами};

\node[flowaction, below=of start] (count)
{Определите неизвестные коэффициенты\\и необходимое число независимых уравнений};

\node[flowaction, below=of count] (points)
{Найдите узловые точки графика\\с точными и удобными координатами};

\node[flowaction, below=of points] (choose)
{Выберите точки с небольшими координатами\\и подставьте \(x\) и \(y=f(x)\) в формулу};

\node[flowdecision, below=13mm of choose] (info)
{Полученное уравнение содержит нужные неизвестные?};

\node[flowaction, right=30mm of info] (other)
{Возьмите другую\\узловую точку};

\node[flowaction, below=15mm of info] (system)
{Составьте систему уравнений\\для коэффициентов};

\node[flowaction, below=of system] (solve)
{Решите систему и восстановите\\конкретную формулу функции};

\node[flowaction, below=of solve] (check)
{Проверьте формулу подстановкой\\координат исходных точек};

\node[flowaction, below=of check] (target)
{Выполните требование задачи:\\найдите \(f(c)\) или решите \(f(x)=c\)};

\node[flowstart, below=of target] (finish)
{Ответ};

\draw[flowarrow] (start) -- (count);
\draw[flowarrow] (count) -- (points);
\draw[flowarrow] (points) -- (choose);
\draw[flowarrow] (choose) -- (info);

\draw[flowarrow]
(info.east) --
node[above,font=\small]{нет}
(other.west);

\draw[flowarrow]
(other.north)
|- ([xshift=12mm]points.east)
|- (points.east);

\draw[flowarrow]
(info.south) --
node[right,font=\small]{да}
(system.north);

\draw[flowarrow] (system) -- (solve);
\draw[flowarrow] (solve) -- (check);
\draw[flowarrow] (check) -- (target);
\draw[flowarrow] (target) -- (finish);

\end{tikzpicture}
\end{center}

\clearpage

% =================================================
\section{7. Типичные ошибки}

\begin{enumerate}

\item
\important{Неверный сдвиг запятой.}

При переходе

\[
3{,}2\cdot10^{-8}
=
32\cdot10^{-9}
\]

показатель уменьшается на \(1\).

\item
\important{Сокращение слагаемых вместо множителей.}

Из

\[
\frac{9}{3^\alpha}
\]

полезно получить

\[
\frac{3^2}{3^\alpha}
=
3^{2-\alpha}.
\]

\item
\important{Вычитание показателей при сложении или вычитании степеней.}

Формула

\[
\frac{a^m}{a^n}=a^{m-n}
\]

относится к делению.

Из

\[
a^m-a^n
\]

нельзя сразу получить

\[
a^{m-n}.
\]

\item
\important{Ошибочное расщепление степени суммы или разности.}

Верно:

\[
a^{m-n}
=
\frac{a^m}{a^n}.
\]

Ошибочно:

\[
a^{m-n}=a^m-a^n.
\]

\item
\important{Автоматический выбор точки \(x=0\).}

Точка может дать тождество вида

\[
1=1
\]

и не сообщить ничего об искомом коэффициенте.

\item
\important{Перепутаны координаты точки.}

Для точки

\[
(x_0;y_0)
\]

первая координата --- аргумент \(x_0\), вторая --- значение функции:

\[
y_0=f(x_0).
\]

\item
\important{Подмена \(f(c)\) и условия \(f(x)=c\).}

\[
f(6)
\]

означает подстановку \(x=6\).

Условие

\[
f(x)=6
\]

означает уравнение, из которого требуется найти \(x\).

\end{enumerate}

% =================================================
\section{8. Короткий чек-лист перед ответом}

Перед завершением задачи проверьте:

\begin{itemize}
    \item одинаковые ли основания у степеней;
    \item действительно ли выполняется умножение или деление там, где Вы работаете с показателями;
    \item верно ли преобразована десятичная дробь;
    \item правильно ли переведены слова «увеличился в \(k\) раз» и «уменьшился в \(k\) раз»;
    \item все ли подставленные координаты принадлежат выбранным точкам;
    \item достаточно ли выбранных точек для нахождения всех коэффициентов;
    \item не исчез ли неизвестный коэффициент после подстановки;
    \item проверена ли найденная формула хотя бы на одной исходной точке;
    \item правильно ли прочитано требование: \(f(c)\) или \(f(x)=c\).
\end{itemize}

% =================================================
% ТРЕНИРОВОЧНЫЙ БЛОК
% =================================================
\clearpage

\section*{Тренировочный блок — Путь А}

\textbf{5 заданий.}
Решения оформляйте последовательно. В ответ запишите итоговое значение.

\begin{enumerate}[label=\textbf{\arabic*.}]

\item
Упростите выражение
\[
\frac{2^{x+5}}{2^{x-1}}.
\]

\item
Решите уравнение
\[
4^{x-2}=0{,}25.
\]

\item
Для некоторого процесса выполняется закон
\[
pV^\alpha=\text{const}.
\]

При переходе из первого состояния во второе объём уменьшился в \(2\) раза, а давление увеличилось в \(8\) раз.

Найдите \(\alpha\).

\item
График функции
\[
f(x)=a^x+b
\]
проходит через точки
\[
A(0;-3),
\qquad
B(1;-1).
\]

Найдите
\[
f(4).
\]

\item
График функции
\[
f(x)=a^x+b
\]
проходит через точки
\[
A(-1;1),
\qquad
B(1;4).
\]

Найдите значение \(x\), при котором
\[
f(x)=7.
\]

\end{enumerate}

% =================================================
% ОТВЕТЫ
% =================================================
\clearpage

\section*{Ответы}

\renewcommand{\arraystretch}{1.45}

\begin{table}[ht]
\centering
\begin{tabularx}{0.82\linewidth}{@{}>{\bfseries}c X@{}}
\toprule
Задание & Ответ \\
\midrule
1 & \(64\) \\
2 & \(1\) \\
3 & \(3\) \\
4 & \(29\) \\
5 & \(2\) \\
\bottomrule
\end{tabularx}
\end{table}

\vfill

\begin{center}
\color{softgray}
\small
Лёвин Артём Александрович
\end{center}

\end{document}